\documentclass[10pt,a4paper]{amsart}
\usepackage[margin=19mm]{geometry}
\usepackage{amsmath,amssymb,mathtools}
\usepackage[scr=rsfs,cal=euler]{mathalfa}
\usepackage{xfrac}
\usepackage[unicode,hidelinks,bookmarks=false]{hyperref}
\newcommand{\ie}{i.e.\ }
\newcommand{\cf}{cf.\ }
\newcommand{\resp}{respectively\ }
\newcommand{\Subsection}{\subsection}
\providecommand{\nobreakdash}{\nobreak\hbox{-}}
\setlength{\emergencystretch}{2em}
\multlinegap0pt
\newcommand{\Ric}{\mathrm{Ric}}
\newcommand{\E}{\mathbb{E}}
\newcommand{\Var}{\mathrm{Var}}
\newcommand{\dd}{\,\mathrm{d}}
%\newcommand{\de}{d_e}
\newcommand{\de}{\tilde d}

\theoremstyle{plain}
\newtheorem{theorem}{Theorem}
\newtheorem{lemma}[theorem]{Lemma}
\newtheorem{proposition}[theorem]{Proposition}
\theoremstyle{remark}
\newtheorem{remark}[theorem]{Remark}
\title[Variance-Refined In-Diameter Lower Bound for the First Dirichlet Eigenvalue]{Variance-Refined In-Diameter Lower Bound for the First Dirichlet Eigenvalue: the variance-refined in-diameter bound (Theorem 1.1)}
\author{Thomas Sch\"urmann}
\date{Source: arXiv:2512.21517v1 (CC BY 4.0)}
\begin{document}
\maketitle

\noindent\emph{Source and licence.} Variance-Refined In-Diameter Lower Bound for the First Dirichlet Eigenvalue. Version arXiv:2512.21517v1 (CC BY 4.0), distributed under the Creative Commons Attribution 4.0 International licence, http://creativecommons.org/licenses/by/4.0/, as recorded in the arXiv licence metadata for this exact version and shown on the abstract page https://arxiv.org/abs/2512.21517v1. Full licence notice: licenses/arxiv-2512.21517-CC-BY-4.0.txt.\par\smallskip

\noindent\textit{Editorial scope.} Counted unit: the article's main theorem, printed as Theorem 1.1 (source0.tex lines 74--97, source label \texttt{thm:main}), the variance-refined in-diameter lower bound $\lambda\ge\bigl((\alpha+D)+\sqrt{(\alpha+D)^2+V\alpha^2}\bigr)/2$ together with its strict improvement over Ling's estimate for $K>0$, delivered with the article's complete original proof (source0.tex lines 508--538, one \texttt{proof} environment of 31 lines carrying the optional argument ``Proof of Theorem~\ref{thm:main}''; the article states the theorem in its first section and prints the proof at the end of its section 5, so the proof is detached from the statement only by the article's own arrangement, and both spans are delivered here). No mathematics is written, repaired or re-derived: statement and proof are the article's own text line for line, in the article's own notation ($\alpha$, $D$, $V$, $\delta$, $\de$, $\xi$, $z$, $\lambda$, $\zeta(3)$), and no line of the counted statement or of the counted proof is omitted, substituted or re-flowed. Required context retained uncounted: the article's standing geometric assumptions and the definition of the first Dirichlet eigenvalue (lines 42--49); the article's own statement of Ling's estimate, which is the bound being refined (lines 58--72; its only citation, to \texttt{Ling2006}, is transcribed in the wrapper's bibliography); the article's Lemma 2.1 on the auxiliary function $\xi$ (lines 109--122); its Lemma 2.2, Ling's integral inequality (lines 151--165); its Proposition 3.1, the strong-convexity variance improvement for $x^{-1/2}$ (lines 191--208); its Lemma 3.2, the mean and variance of $z$ (lines 249--266); its Proposition 3.3, the variance-refined integral estimate (lines 279--290); its Lemma 4.1, the second moment of $\xi$ and the closed form of $V$ (lines 304--318) together with the alternative form of $\xi$ that it uses (lines 324--327); its Remark 4.2 recording $V>0$ (lines 386--391); its Lemma 5.1, the one-root lower bound (lines 397--407); and its Lemma 5.3, the range of $\delta$ (lines 493--499). Editorial formatting only, in addition: an AMS \texttt{amsart} preamble that restates the article's own macro definitions (\texttt{\textbackslash Ric}, \texttt{\textbackslash E}, \texttt{\textbackslash Var}, \texttt{\textbackslash dd}, \texttt{\textbackslash de}) and the theorem environments the retained lines use; sequential numbering of the retained environments, so that they are numbered in order of appearance, whereas the article prints them as Theorem 1.1, Lemma 2.1, Lemma 2.2, Proposition 3.1, Lemma 3.2, Proposition 3.3, Lemma 4.1, Remark 4.2, Lemma 5.1 and Lemma 5.3; the article's section-relative equation numbering is likewise not reproduced, so displayed equations inside the retained spans are numbered consecutively instead, and every cross-reference inside the retained text resolves to the wrapper's numbering. The article's class file and its package list are not shipped; the wrapper is the pinned AMS runtime. Exactly one bibliography entry is delivered, the cited key \texttt{Ling2006}, transcribed character for character from the article's own \texttt{thebibliography} block (source0.tex lines 693--717); no other reference is cited inside a retained span. No asset-inclusion command occurs inside any retained line, so the wrapper carries no asset.

\section*{Standing assumptions and the estimate being refined}

Let $(M,g)$ be a compact $n$-dimensional Riemannian manifold with nonempty boundary $\partial M$.
Assume that the Ricci curvature satisfies
\begin{equation}
\label{eq:Ric-lower}
\Ric(M)\ge (n-1)K
\end{equation}
for some constant $K>0$, and that the mean curvature of $\partial M$ with respect to the outward unit normal is nonnegative.
Let $\lambda$ denote the first Dirichlet eigenvalue of the Laplacian on $M$.


In \cite{Ling2006}, Ling proved a unified in-diameter estimate.
Writing
\begin{equation}
\de = 2\sup_{x\in M}\mathrm{dist}(x,\partial M),
\end{equation}
Ling's main theorem yields
\begin{equation}
\lambda \ge \frac{(n-1)K}{2}+\frac{\pi^2}{\de^{\,2}}.
\end{equation}
The proof is based on a refined gradient comparison and introduces an auxiliary function $\xi$ on $[-\pi/2,\pi/2]$.
After reducing the eigenvalue estimate to a one-dimensional integral inequality, Ling applies H\"older's inequality to replace the nonconstant comparison function $z(t)$ by its mean.
This step discards information about the oscillation of $z$.

The aim of this note is to retain this oscillation quantitatively via a variance refinement.
We obtain an explicit strengthening of Ling's bound in closed form.


\section*{Retained context: the auxiliary function and Ling's integral inequality}

\begin{lemma}[The auxiliary function $\xi$]
\label{lem:xi}
Define $\xi:[-\pi/2,\pi/2]\to\mathbb{R}$ by
\begin{equation}
\label{eq:xi-def}
\xi(t)=\frac{\cos^2 t + 2t\sin t\cos t + t^2 - \pi^2/4}{\cos^2 t}.
\end{equation}
Then $\xi$ is smooth and even on $[-\pi/2,\pi/2]$, satisfies $\xi(\pm \pi/2)=0$, and
\begin{equation}
\label{eq:xi-mean}
\int_{0}^{\pi/2}\xi(t)\,\dd t = -\frac{\pi}{2}.
\end{equation}
Moreover $\xi(t)\le 0$ for $t\in[0,\pi/2]$.
\end{lemma}


\begin{lemma}[Ling's integral inequality]
\label{lem:ling-integral}
Let $\lambda$ be the first Dirichlet eigenvalue, set $\alpha=(n-1)K/2$ and $\delta=\alpha/\lambda$.
Define
\begin{equation}
\label{eq:z-def}
z(t)=1+\delta\,\xi(t),
\qquad t\in[0,\pi/2].
\end{equation}
Then
\begin{equation}
\label{eq:ling-star}
\sqrt{\lambda}\,\frac{\de}{2}\ \ge\ \int_{0}^{\pi/2}\frac{\dd t}{\sqrt{z(t)}}.
\end{equation}
\end{lemma}


\section*{Retained context: the variance refinement}

\begin{proposition}[Variance improvement for $x^{-1/2}$]
\label{prop:variance}
Let $z:[0,\pi/2]\to(0,1]$ be measurable and set
\begin{equation}
\mu = \frac{2}{\pi}\int_{0}^{\pi/2} z(t)\,\dd t.
\end{equation}
Then
\begin{equation}
\label{eq:variance-ineq}
\frac{2}{\pi}\int_{0}^{\pi/2}\frac{\dd t}{\sqrt{z(t)}}
\ \ge\
\frac{1}{\sqrt{\mu}}+\frac{3}{8}\,\Var(z),
\end{equation}
where
\begin{equation}
\Var(z)=\frac{2}{\pi}\int_{0}^{\pi/2}\bigl(z(t)-\mu\bigr)^2\,\dd t.
\end{equation}
\end{proposition}


\begin{lemma}[Mean and variance of $z(t)=1+\delta\xi(t)$]
\label{lem:meanvar-z}
Let $z(t)=1+\delta\xi(t)$ on $[0,\pi/2]$, where $\xi$ is given by \eqref{eq:xi-def}.
Then
\begin{equation}
\label{eq:mu-1-delta}
\mu = \frac{2}{\pi}\int_0^{\pi/2} z(t)\,\dd t = 1-\delta,
\end{equation}
and
\begin{equation}
\label{eq:var-z}
\Var(z)=\delta^2\,\Var(\xi),
\qquad
\Var(\xi)=\E[\xi^2]-1,
\qquad
\E[\xi^2]=\frac{2}{\pi}\int_0^{\pi/2}\xi(t)^2\,\dd t.
\end{equation}
\end{lemma}


\begin{proposition}[Variance-refined integral estimate]
\label{prop:refined-integral}
Let $z(t)=1+\delta\xi(t)$ as in \eqref{eq:z-def}.
Then
\begin{equation}
\label{eq:integral-refined}
\int_{0}^{\pi/2}\frac{\dd t}{\sqrt{z(t)}}
\ \ge\
\frac{\pi}{2}\left(\frac{1}{\sqrt{1-\delta}}+\frac{3}{8}\Var(\xi)\,\delta^2\right).
\end{equation}
In particular, if $\Var(\xi)>0$ and $\delta>0$, the right-hand side is strictly larger than $\frac{\pi}{2}(1-\delta)^{-1/2}$.
\end{proposition}


\section*{Retained context: the constant $V$ and the one-root majorization}

\begin{lemma}[Second moment of $\xi$]
\label{lem:xi-secondmoment}
For $\xi$ defined by \eqref{eq:xi-def},
\begin{equation}
\label{eq:xi-square-integral}
\int_{0}^{\pi/2}\xi(t)^2\,\dd t
=
\pi\left(2\zeta(3)-\frac{\pi^2+1}{6}\right).
\end{equation}
Consequently,
\begin{equation}
\label{eq:V-closedform}
V = 4\,\zeta(3)-\frac{1}{3}(\pi^2+4).
\end{equation}
\end{lemma}


\begin{equation}
\label{eq:xi-alt}
\xi(t) = 1 + 2t\tan t + \left(t^2-\frac{\pi^2}{4}\right)\sec^2 t.
\end{equation}


\begin{remark}\label{rem:V-positive}
The constant $V$ is a variance, hence nonnegative by definition.
Moreover $V>0$ because $\xi$ is not (a.e.) constant: using \eqref{eq:xi-alt} one has
\(\xi(0)=1-\pi^2/4<0\) (since $\pi>2$), while $\xi(\pi/2)=0$ by Lemma~\ref{lem:xi}.
For reference, $V\approx 0.1850261456$.
\end{remark}


\begin{lemma}[One-root lower bound]
\label{lem:one-root}
Let $V>0$ be as in \eqref{eq:V-closedform}.
For every $\delta\in[0,1/2]$ one has
\begin{equation}
\label{eq:one-root}
\frac{1}{\sqrt{1-\delta}}+\frac{3}{8}V\,\delta^2
\ \ge\
\frac{1}{\sqrt{1-\delta-\frac{V}{4}\delta^2}}.
\end{equation}
\end{lemma}


\begin{lemma}[Range of $\delta$]
\label{lem:delta-range}
Under the standing geometric assumptions,
\begin{equation}
0\le \delta=\frac{\alpha}{\lambda}\le \frac{n-1}{2n}<\frac{1}{2}.
\end{equation}
\end{lemma}


\section{The counted unit: Theorem 1.1 and its complete original proof}

\begin{theorem}[Variance-refined in-diameter bound]
\label{thm:main}
Let $(M,g)$ satisfy the assumptions above and let $\lambda$ be the first Dirichlet eigenvalue.
Set
\begin{equation}
\alpha = \frac{(n-1)K}{2},
\qquad
D=\frac{\pi^2}{\de^{\,2}},
\qquad
V = 4\zeta(3)-\frac{1}{3}(\pi^2+4).
\end{equation}
Then
\begin{equation}
\label{eq:mainbound}
\lambda\ \ge\
\frac{(\alpha+D)+\sqrt{(\alpha+D)^2+V\,\alpha^2}}{2}.
\end{equation}
In particular, since $V=\Var(\xi)>0$ (see Remark~\ref{rem:V-positive}), one has the strict improvement
\begin{equation}
\lambda\ >\ \alpha + D
\qquad
\text{whenever }K>0.
\end{equation}
\end{theorem}

\begin{proof}[Proof of Theorem~\ref{thm:main}]
Set $\delta=\alpha/\lambda$.
By Lemma~\ref{lem:ling-integral}, Proposition~\ref{prop:refined-integral}, Lemma~\ref{lem:xi-secondmoment}, and Lemmas~\ref{lem:one-root}--\ref{lem:delta-range}, we have
\begin{equation}
\sqrt{\lambda}\,\frac{\de}{2}
\ \ge\
\int_{0}^{\pi/2}\frac{\dd t}{\sqrt{z(t)}}
\ \ge\
\frac{\pi}{2}\cdot\frac{1}{\sqrt{1-\delta-\frac{V}{4}\delta^2}}.
\end{equation}
Squaring and writing $D=\pi^2/\de^{\,2}$ yields
\begin{equation}
\label{eq:lambda-rational}
\lambda \ge \frac{D}{1-\delta-\frac{V}{4}\delta^2}.
\end{equation}
Substitute $\delta=\alpha/\lambda$ into \eqref{eq:lambda-rational} and clear denominators:
\begin{equation}
\lambda\left(1-\frac{\alpha}{\lambda}-\frac{V}{4}\frac{\alpha^2}{\lambda^2}\right)\ge D.
\end{equation}
Equivalently,
\begin{equation}
\lambda^2-(\alpha+D)\lambda-\frac{V}{4}\alpha^2\ge 0.
\end{equation}
Since $\lambda>0$, this quadratic inequality implies
\begin{equation}
\lambda\ge \frac{(\alpha+D)+\sqrt{(\alpha+D)^2+V\,\alpha^2}}{2},
\end{equation}
which is \eqref{eq:mainbound}.

Finally, $V>0$ by \eqref{eq:V-closedform}, so the square-root term is strictly larger than $\alpha+D$ whenever $\alpha>0$, proving the strict improvement when $K>0$.
\end{proof}

\begin{thebibliography}{99}
\bibitem[Ling2006]{Ling2006} J.~Ling, \emph{A lower bound of the first Dirichlet eigenvalue of a compact manifold with positive Ricci curvature}, International Journal of Mathematics \textbf{17} (2006), no.~5, 605--617.
\end{thebibliography}
\end{document}
